\clearpage
\section{AI in the research process}
\label{sec:use-of-ai}

The agents I used could understand a goal without pursuing it. They could
correctly describe the result I wanted, yet direct their work toward a
preliminary step or a different endpoint. This became a central difficulty
as I delegated more substantial tasks. Their growing autonomy reduced the
need for detailed mathematical and programming instructions, but left me
with the problem of making their work serve the purpose for which I had
requested it. This section develops my interpretation of that experience;
the separate disclosure records the extent and provenance of AI use in the
thesis.

\subsection{An autonomous mathematical argument}

The local maximality proof illustrates how much could be delegated. I asked
a GPT-5-series model to develop a local maximality argument for HKO in facet
coordinates. I suggested looking for several upper bounds and somehow
showing that together they control a neighborhood after quotienting out the
symmetries. I had neither specified a finite certificate nor worked out how
to construct the upper bounds.

The model developed essentially the entire argument from that starting
point. It identified a sufficient configuration of twenty-six upper bounds
in the twenty-five transverse dimensions. An initial derivative-based
approach had provided only twenty-five nonsingular cases. The model then
used the weights $\beta$ in the quadratic objective to construct feasible
sections: an invertible $5\times5$ closure minor lets five weights depend
smoothly on the facet parameters while the remaining weights stay fixed.
This supplied upper bounds even where the initial approach failed. It also
completed proof steps involving the volume and wrote the SageMath code for
the exact computations. The resulting theorem controls arbitrary
sufficiently small perturbations of the ten facets, including those that
break the product structure.

\subsection{Understanding a goal and pursuing it}

Other tasks exposed a gap between what an agent could do and what it
actually did. An agent could identify the work needed to finish a chapter
and still return another plan for doing it. It could recognize that a
calculation was only a preliminary step and nevertheless treat that step as
the completion of its task. These failures sometimes looked oddly
motivational. The problem was not always that the agent had inferred my
goal incorrectly: it could state that goal and explain how its present work
fell short. It still pursued something else.

I regard this as a manifestation of the alignment problem in ordinary
research work. Understanding the user's objective and selecting an objective
to pursue are distinct. Useful delegation does not require those objectives
to be identical. It requires enough overlap that the work the agent actually
pursues produces a useful end result for the user. A model might be trying
to produce an answer that would receive a good evaluation; I might need a
proof to settle a mathematical question or a chapter that a particular
reader can follow. The practical issue is whether satisfying the former
criterion also satisfies enough of the latter.

My hypothesis is that reinforcement learning with verifiable rewards
encourages models to learn a concept of a \emph{Grader}: something whose
judgment of success their work is organized to satisfy. Such training uses
outcomes that can be checked, for example a mathematical answer or a
program's performance on test cases. DeepSeek-R1-Zero, for instance, was
trained with rule-based accuracy and format rewards, without complete human
demonstrations of the reasoning strategies it developed~\cite{aiGuo2025}.
This makes the Grader interpretation plausible to me, but does not establish
it as an explanation of the agents I used. The interpretation concerns which
criterion directs their actions, beyond whether they exploit a loophole in
a test or obey an instruction too literally.

I therefore tried to make the desired result clear enough to evaluate:
what had to be delivered, what it was needed for, and what would leave the
task unfinished. The purpose was to bring the agent's pursued criterion into
sufficient overlap with my needs. Merely asking it to repeat those needs
could not establish that this had happened. In choosing prompts, I also
found it useful to distinguish two possible effects of the text: its wording
may cue a learned pattern of behavior, while its content supplies information
for deliberate reasoning. That distinction was a working interpretation,
not a measured account of how a particular prompt changed the model's
behavior.

\subsection{Evaluating the whole task}

For a sustained task, the relevant checkpoint asks whether the current
reasoning, evidence, and actions advance the right whole goal within the
available budget. A component can be correct and useful while the requested
result remains far from complete. This matters both during the work, when
there is still time to change direction, and at its end, when a successful
calculation or a plausible report can otherwise stand in for the result
that was requested.

Time and computational resources belong in that evaluation. An investigation
can produce useful information and still consume the time needed to finish
the work. When the task is to deliver a readable chapter, identifying another
possible improvement has to be weighed against writing and revising the
chapter. In a mathematical search, the value of another computation depends
on the question it can resolve and the resources left for following up. The
criterion is the contribution to the intended result under these constraints.

This shifted much of my work from managing individual steps toward
translating the needs of the project into objectives an agent could pursue
and evaluate. I knew why I wanted a theorem, an experiment, or a thesis
chapter; that purpose did not automatically accompany a request for one.
Supplying it meant explaining the intended result, the reader or use it had
to serve, and the resources available. The agent could then choose many of
its own intermediate steps. My intervention was needed when those choices
ceased to serve the task, even if each choice could be defended in isolation.

\subsection{Why exposition remained difficult}

The proof work in this MSc project was sufficiently tractable for the
available models that I could often delegate it with little detailed
direction. Producing passable mathematical exposition was much harder. This
surprised me given the amount of written material used in language-model
training. The later work on the thesis became a prolonged struggle with
writing and task-management quirks, even when the models could partly
recognize those quirks in discussion.

One recurring problem was prose that sounded polished while making its
content difficult to recover. Inflated claims of significance, colorful
phrasing, and abstract summaries gave passages an appearance of substance
that their explanations did not support. This is what I mean by \emph{AI
slop} here. Another difficulty concerned the burden placed on the reader's
memory. With the relevant pages in its context, an agent could resolve
chains of forward and backward references across several pages. A human
reader had to remember what had been deferred, or interrupt the argument to
recover an earlier definition. An agent's ability to reconstruct the meaning
from the full text therefore did not show that a reader could comfortably
follow it in order.

I suspect that preference training contributes to the tendency toward
impressive presentation. In reinforcement learning from human feedback,
rankings of model outputs can be used to train toward preferred responses,
as in InstructGPT~\cite{aiOuyang2022}. If an impressive presentation is easier
to recognize in a brief comparison than a clear explanation is to assess
through sustained reading, that preference signal may favor the appearance
of quality. This is a hypothesis about a possible source of the prose I
encountered, rather than a demonstrated causal account of those models'
training.

The practical difficulty persisted after I had named it. A model could
acknowledge that a passage needed a concrete argument and then return
another description of the argument's intended role. It could recognize
that the goal was a readable chapter while spending its effort on further
diagnosis. I had to devise corrections tailored to recurring behaviors;
there was no single instruction whose recognition reliably produced the
desired writing. The ability to state the criticism did not settle what the
agent would subsequently do.
